% Fuente: Control 1, P1.
El problema de coordinación y ruteo de equipos de mantenimiento es fundamental para reducir los costos de operación en parques eólicos. En este contexto, la empresa \textit{WindCore Maintenance} desea planificar el mantenimiento de un conjunto de turbinas eólicas distribuidas geográficamente, coordinando simultáneamente la asignación de equipos técnicos y sus rutas de desplazamiento.

Sea $\mathcal{T}=\{1,\dots,T\}$ el conjunto de períodos de planificación, $\mathcal{I}=\{1,\dots,I\}$ el conjunto de turbinas eólicas y $\mathcal{J}=\{1,\dots,J\}$ el conjunto de equipos técnicos disponibles. Todas las turbinas se encuentran en ubicaciones $i\in\mathcal{I}$, y existe un depósito central, denotado por $\delta$, desde donde los equipos técnicos inician y finalizan sus rutas en cada período. Sea $\mathcal{V}=\mathcal{I}\cup\{\delta\}$ el conjunto de nodos del problema. El costo de traslado entre dos ubicaciones $i_1,i_2\in\mathcal{V}$ está dado por $C_{i_1i_2}^{tl}$. Este costo es independiente del equipo técnico utilizado.

A partir de información entregada por sensores, se conoce de antemano la condición de cada turbina en cada período. Esta información se incorpora directamente mediante el parámetro $\bar{C}_{it}$, que representa el costo de realizar mantenimiento a la turbina $i$ en el período $t$. Cada equipo técnico $j\in\mathcal{J}$ posee un nivel de experiencia que determina tanto su capacidad para realizar el mantenimiento como el tiempo requerido para ello. En particular, $\bar{o}_{ijt}$ representa el tiempo requerido por el equipo $j$ para realizar mantenimiento a la turbina $i$ en el período $t$, mientras que $a_{ijt}$ es un parámetro binario que toma valor $1$ si el equipo $j$ está capacitado para realizar dicho mantenimiento, y $0$ en caso contrario. Cada turbina debe recibir exactamente un mantenimiento durante el horizonte de planificación.

En cada período $t$, un equipo técnico puede ser desplegado desde el depósito para visitar una o más turbinas, realizar secuencialmente los mantenimientos programados y retornar al depósito. Por lo tanto, cada \textit{tour} debe completarse dentro de un mismo período $t$. Alternativamente, el equipo puede no ser desplegado durante dicho período. La duración total del \textit{tour}, incluyendo tiempos de traslado y mantenimiento, no puede superar la jornada laboral disponible $\Delta_T$. Así, cada equipo debe finalizar su ruta y regresar al depósito dentro del tiempo disponible. El tiempo de traslado entre dos ubicaciones $i_1,i_2\in\mathcal{V}$ se denota por $\tau_{i_1i_2}$.


Adicionalmente, se incorporan vehículos eléctricos para el desplazamiento de los equipos técnicos. Estos vehículos poseen una autonomía limitada, por lo que la distancia total recorrida por cada equipo durante un período de trabajo no puede superar un valor máximo $E$. La distancia entre dos ubicaciones $i_1,i_2\in\mathcal{V}$ se denota por $d_{i_1i_2}$. Asimismo, debido a restricciones de capacidad para transportar repuestos, cada vehículo puede atender como máximo $M$ turbinas en un mismo \textit{tour}.


\textit{Hint:} \textit{Para formular las restricciones de eliminación de subtours, puede interpretar cada ruta como un problema de abastecimiento desde el depósito. En esta interpretación, cada turbina consume una unidad de capacidad solo si se planifica mantenimiento en ella, y cada vehículo dispone de una capacidad máxima de $M$ unidades. Podría ser útil introducir  variables auxiliares de flujo.}

 
\begin{enumerate}
    \item Plantee un modelo MILP para determinar la planificación de mantenimiento y las rutas de los equipos técnicos, minimizando el costo total de mantenimiento y traslado, sujeto a las restricciones descritas anteriormente. Considere que todas las turbinas eólicas deben recibir mantenimiento exactamente una vez durante el periodo de planificación.  
    
    \item Debido a los efectos del calentamiento global, en algunos períodos las condiciones climáticas pueden impedir que los equipos de mantenimiento realicen las actividades programadas. Para anticipar esta situación, se dispone de un sistema de monitoreo meteorológico de última generación, cuya información se representa mediante el parámetro binario $W_t$. Este parámetro toma el valor $1$ si las condiciones climáticas permiten realizar mantenimiento en el período $t$, y $0$ en caso contrario. Indique qué modificaciones deben incorporarse al modelo para considerar esta información meteorológica en la planificación de mantenimiento y ruteo, obteniendo un nuevo modelo MILP.
    
    \item Debido a la creación reciente de un sindicato, originada por una alta carga laboral en algunos equipos de mantenimiento, se estableció un acuerdo que exige balancear la carga de trabajo entre equipos. En particular, la diferencia entre las horas totales trabajadas por cualquier par de equipos durante el horizonte de planificación no puede superar un valor máximo $\gamma$. Indique las restricciones necesarias para incorporar esta condición al modelo, manteniendo una formulación MILP.
\end{enumerate}
